\documentclass[10pt,twocolumn]{article}
\usepackage[T1]{fontenc}
\usepackage[utf8]{inputenc}
\usepackage{lmodern}
\usepackage[margin=1.8cm,top=2.4cm,bottom=2cm,columnsep=0.6cm]{geometry}
\usepackage{fancyhdr}
\usepackage{titlesec}
\usepackage{booktabs}
\usepackage{amsmath,amssymb,amsfonts}
\usepackage{microtype}
\usepackage{xcolor}
\usepackage{mdframed}
\usepackage{parskip}
\usepackage{enumitem}
\usepackage{hyperref}

\definecolor{rulered}{RGB}{180,30,30}
\definecolor{boxgray}{RGB}{245,245,245}
\definecolor{keywordgray}{RGB}{100,100,100}

\pagestyle{fancy}
\fancyhf{}
\fancyhead[L]{\textsc{\small Claude Mag}}
\fancyhead[C]{\small\textit{Machine Learning Research Digest}}
\fancyhead[R]{\small 2026-10-07}
\fancyfoot[C]{\small\thepage}
\renewcommand{\headrulewidth}{0.8pt}

\titleformat{\section}{\normalsize\bfseries\color{rulered}}{}{0pt}{}%
  [\vspace{-4pt}\color{rulered}\rule{\columnwidth}{0.4pt}]
\titlespacing{\section}{0pt}{8pt}{4pt}

\hypersetup{colorlinks=true,urlcolor=rulered,linkcolor=black}

\begin{document}
\setlength{\parindent}{0pt}
\setlength{\parskip}{4pt}

{\small\color{keywordgray}\textbf{Keywords:}
  world models \textbullet{}
  hierarchical planning \textbullet{}
  JEPA \textbullet{}
  latent representation}\par\vspace{4pt}

{\LARGE\bfseries\raggedright
  H-JEPA: End-to-End Learning of Hierarchical World Models for
  Visual Planning}\par\vspace{4pt}

{\large\itshape\raggedright
  Stacking Action-Conditioned JEPAs Hierarchically Raises Visual AntMaze
  Success from 18\% to 73\% by Letting Each Level Predict Farther Ahead in
  Its Own Latent Space}\par\vspace{6pt}

{\small\textbf{By} Wancong Zhang, Basile Terver, Michael Rabbat, Yann LeCun
  \& Randall Balestriero (Meta AI; New York University)}\par\vspace{2pt}

{\small\color{keywordgray}arXiv:2610.06805 \textbullet{} October 2026}
\par\vspace{2pt}\noindent{\color{rulered}\rule{\columnwidth}{0.6pt}}\vspace{6pt}

Flat world models trained at a single temporal resolution conflate the
timescales of muscle-level control, motion planning, and goal-directed
navigation, forcing planners to reason about irrelevant visual detail.
H-JEPA stacks action-conditioned JEPAs so that higher levels automatically
discard fast-varying sensory noise and retain the slow, task-relevant state
that long-horizon planning depends on---without any manual specification
of sub-goals or timescale boundaries.

\begin{mdframed}[backgroundcolor=boxgray,linecolor=rulered,linewidth=1pt,
    innerleftmargin=6pt,innerrightmargin=6pt,
    innertopmargin=6pt,innerbottommargin=6pt]
\textbf{\small AT A GLANCE}\par\vspace{4pt}
\begin{description}[leftmargin=0pt,labelsep=4pt,itemsep=2pt,parsep=0pt,
    font=\small\bfseries]
  \item[Problem:] Single-level world models force planners to reason at one
    temporal resolution, wasting capacity on irrelevant low-level detail
    during long-horizon planning.
  \item[Why it matters:] Robotics, navigation, and embodied AI require
    simultaneous reasoning across sub-second control and minute-scale goals.
  \item[Method:] Stack $L$ action-conditioned JEPAs; each level predicts
    $k$ steps ahead in the latent of the level below; train all levels
    jointly end-to-end with a VICReg-style regularizer.
  \item[Result:] A 3-level hierarchy raises Visual AntMaze success from
    \textbf{18\% to 73\%} while reducing planner compute.
  \item[Evidence:] Visual AntMaze, robotic manipulation, and real-robot
    planning on the DROID dataset.
\end{description}
\end{mdframed}

\section{The Big Picture}

Long-horizon planning requires simultaneous reasoning across radically
different timescales: sub-second muscle control, second-scale motion planning,
and minute-scale goal navigation. A flat world model trained at a single
temporal resolution cannot separate these concerns---it must predict every
future state at the same resolution, forcing the planner to process
irrelevant low-level visual detail even when making high-level route decisions.

H-JEPA (Hierarchical Joint Embedding Predictive Architecture) resolves this
by stacking multiple levels of action-conditioned JEPAs, each operating at
a coarser timescale than the one below. The hierarchy is not hand-designed:
timescale separation emerges automatically from end-to-end training when the
world contains factors that evolve at different speeds.

\section{Why Existing Approaches Fall Short}

\textbf{Flat world models are computationally wasteful.} A single-level model
predicts every future state at native resolution, burdening the planner with
fine-grained sensory detail even during high-level decision making.

\textbf{Hand-designed hierarchies require manual sub-goals.} Options
frameworks and HIRO-style methods pre-specify the timescale structure and
require human-designed sub-goal spaces, sacrificing end-to-end trainability.

\textbf{Generative hierarchical models are fragile.} Pixel-prediction
hierarchies accumulate errors across levels and are sensitive to the choice
of reconstruction loss.

\textbf{Flat JEPA models remain at a single timescale.} V-JEPA and related
work learn expressive single-level latent predictors; extending them to
multiple timescales in an end-to-end trainable manner is non-trivial.

\section{Core Mechanism}

H-JEPA stacks $L$ levels of action-conditioned JEPAs:

\begin{itemize}[itemsep=2pt,parsep=0pt]
  \item \textbf{Level 1} operates at the native frame rate, predicting the
    latent of the next frame from the current latent and action.
  \item \textbf{Level $\ell > 1$} operates at a $k$-times-coarser timescale.
    It receives the temporally pooled latent from level $\ell{-}1$ and
    predicts that latent $k$ frames ahead.
  \item \textbf{End-to-end training} jointly minimizes a JEPA objective at
    all levels: minimize distance between predicted latent and the
    stop-gradient target from the level below, regularized by a VICReg-style
    term to prevent representational collapse.
\end{itemize}

The key emergent property is \emph{timescale separation}: higher levels learn
to discard fast, unpredictable visual detail and retain slow, task-relevant
state (object positions, goal proximity) without explicit supervision.

\section{Technical Formulation}

Let $z_t^\ell$ denote the latent at level $\ell$ and time $t$. Temporal
pooling $P_k$ averages $k$ consecutive latents. The H-JEPA target at level
$\ell$ is:
\[
z_t^{\ell} = P_k\!\left(z_{t \cdot k}^{\ell-1}, \ldots,
  z_{(t+1) \cdot k - 1}^{\ell-1}\right)
\]

The predictor $f^\ell_\phi$ at level $\ell$ predicts $k$ steps ahead:
\[
\hat{z}_{t+1}^\ell = f^\ell_\phi\!\left(z_t^\ell,\;
  a_{t \cdot k : (t+1) \cdot k}\right)
\]

The per-level JEPA loss is:
\[
\mathcal{L}^\ell = \left\| \hat{z}_{t+1}^\ell
  - \mathrm{sg}(z_{t+1}^\ell) \right\|^2
  + \lambda \cdot \mathcal{L}_\text{VICReg}(z^\ell)
\]

The total objective sums over all levels:
\[
\mathcal{L}_\text{H-JEPA} = \sum_{\ell=1}^{L} \mathcal{L}^\ell
\]

Planning applies a latent MPPI or MCTS planner at the highest level $L$,
scoring rollouts in the level-$L$ latent space---dramatically reducing
compute per planning node relative to planning at level~1.

\section{Learning or Inference Procedure}

\textbf{Training:}
\begin{enumerate}[itemsep=2pt,parsep=0pt]
  \item Collect trajectories $(o_1, a_1, \ldots, o_T, a_T)$.
  \item Encode observations to $z^0 = g_\psi(o)$; propagate up:
    $z^\ell = P_k(z^{\ell-1})$.
  \item Train all levels jointly with $\mathcal{L}_\text{H-JEPA}$.
  \item Optionally add inverse-dynamics heads at each level.
\end{enumerate}

\textbf{Planning at test time:}
\begin{enumerate}[itemsep=2pt,parsep=0pt]
  \item Encode current observation; propagate to level $L$.
  \item Run MPPI or MCTS in the level-$L$ latent space.
  \item Execute the first action; receive next observation; repeat.
\end{enumerate}

\section{What the Guarantee Says}

No formal convergence guarantee is provided. The central empirical claim is
that timescale separation emerges automatically from end-to-end training
when the data contains factors that evolve at different speeds. This is
validated by probing each level's representations for factors of variation
at different timescales. Ablations confirm that a flat JEPA with the same
total parameter count does not match the hierarchical planner.

\section{Experimental Findings}

\begin{itemize}[itemsep=2pt,parsep=0pt]
  \item \textbf{Visual AntMaze:} A 3-level hierarchy raises success rate from
    \textbf{18\% (flat JEPA) to 73\%} while requiring less planner compute.
  \item \textbf{Manipulation:} H-JEPA outperforms flat JEPA on robotic
    pick-and-place and peg-insertion tasks requiring multi-step coordination.
  \item \textbf{DROID (real robot):} With inverse-dynamics supervision,
    H-JEPA improves offline planning fidelity on diverse real-robot
    demonstration data.
  \item \textbf{Timescale analysis:} Level-3 representations encode
    task-relevant state (goal proximity, object position) almost entirely;
    level-1 is dominated by low-level visual texture.
\end{itemize}

\section{Ablations and Interpretation}

\textbf{Number of levels:} Performance improves consistently from 1 to 3
levels; 4 levels adds marginal gain but increases training complexity.

\textbf{Pooling stride $k$:} Smaller strides ($k{=}2$) help tasks with
fast-changing task-relevant state; larger strides ($k{=}8$) help tasks with
very long planning horizons.

\textbf{Timescale separation:} Probing confirms the hierarchy: level~1 is
dominated by visual texture; level~3 is nearly entirely task-state driven.

\textbf{Comparison to hand-designed hierarchies:} H-JEPA matches or exceeds
options-based and HIRO-style methods on standard benchmarks while requiring
no manual specification of sub-goals or timescale boundaries.

\section*{Reference}

Wancong Zhang, Basile Terver, Michael Rabbat, Yann LeCun \& Randall
Balestriero.
\textbf{H-JEPA: End-to-End Learning of Hierarchical World Models for
Visual Planning}.
arXiv:2610.06805, October 2026.\\
\url{https://arxiv.org/abs/2610.06805}

\end{document}
